\clearpage\thispagestyle{empty}
\pdfbookmark[0]{Quellen zu Der Zodiakus der Juno (1805)}{zodiacus1805-sources}
\begin{center}\textbf{Quellen zu Der Zodiakus der Juno (1805)}\end{center}
\noindent Diese KI-Übertragung in heutiges Deutsch folgt Carl Friedrich Gauß’ deutschem Original in den \textit{Werken}, Band~VI (Göttingen, 1874), S.~119--121, einschließlich der vollständigen Tabelle. Der Aufsatz erschien zuerst als Beitrag XX im Märzheft 1805 der von Franz Xaver von Zach herausgegebenen \textit{Monatlichen Correspondenz zur Beförderung der Erd- und Himmels-Kunde}, Band~11, S.~225--228. \href{https://gdz.sub.uni-goettingen.de/id/PPN236007467}{Historischer deutscher Text in den Werken}. \href{https://www.digitale-sammlungen.de/en/view/bsb10538603?page=237}{Erstdruck}.
\par\smallskip\noindent Der Nachtrag bezieht sich auf Gauß’ Aufsatz über die Grenzen der geozentrischen Positionen der Planeten auf S.~106--118 dieses Bandes. Sein \href{https://www.digitale-sammlungen.de/en/view/bsb10538602?page=181}{Erstdruck} erschien im Augustheft 1804, Band~10, Beitrag XIV, S.~173--191.
\par\smallskip\noindent Übertragung und Satz wurden mit KI erstellt.
\clearpage
